% modeling/schwerpunkt_trassierung-DE.tex
\chapter{Schwerpunkt-Trassierung}

\section*{Motivation}

Der klassische Trassierungsentwurf trennt:
\begin{itemize}
\item geometrischen Entwurf (Alignment),
\item Überhöhungsentwurf und
\item dynamische Bewertung.
\end{itemize}

Diese Trennung führt zu suboptimalen Lösungen, weil Geometrie und Dynamik
intrinsisch gekoppelt sind. In der Eisenbahnpraxis zeigt sich diese Kopplung im
Zusammenwirken von Krümmung, Überhöhung, Geschwindigkeit und Fahrzeugreaktion
\cite{Klein1998Trassierung,EN13803_2017,BrustadDalmo2020Review}.

Das Konzept der \emph{Schwerpunkt-Trassierung} will diese Aspekte durch ein
einheitliches Entwurfsprinzip auf Grundlage der Bewegung eines repräsentativen
Punkts verbinden.

Die Vorstellung, Trassierung auf die Bewegung des Fahrzeugschwerpunkts statt
lediglich auf die Gleismittellinie zu beziehen, erscheint auch in
Engineering-Praxis und Patentliteratur \cite{Hasslinger2002Schwerpunkt}.

Dies unterstützt die weiter gefasste AIM-Interpretation eines Alignments als
dynamisch relevante räumliche Zustandstrajektorie und nicht als rein
geometrische Kurve.

\section{Begriffliche Idee}

\begin{definition}[Referenztrajektorie]
Es bezeichne
\[
\gamma_c:I\to\R^2
\]
eine Referenztrajektorie, welche die Bewegung eines charakteristischen Punkts
des Fahrzeugsystems darstellt.
\end{definition}

Dieser Punkt kann interpretiert werden als:
\begin{itemize}
\item Massenschwerpunkt,
\item dynamisch relevanter Bezugspunkt oder
\item abstrakte Entwurfstrajektorie.
\end{itemize}

Der begriffliche Wechsel besteht darin, den Gleisentwurf anhand der Bewegung
eines solchen Bezugspunkts statt allein anhand der Geometrie zu bewerten.

\section{Bezug zur Gleisgeometrie}

Die physische Gleismittellinie \(\gamma(s)\) und der Überhöhungswinkel
\(\phi(s)\) induzieren eine räumliche Konfiguration des Fahrzeugs.

\begin{definition}[Induzierte Referenztrajektorie]
Für einen pose3-Zustand
\[
(\gamma(s),\theta(s),\phi(s))
\]
ist die induzierte Referenztrajektorie
\[
\gamma_c(s)=\gamma(s)+d\,n(s,\phi(s)),
\]
wobei \(d\) ein charakteristischer Versatz und \(n(s,\phi)\) eine von
Orientierung und Überhöhung abhängige Richtung ist.
\end{definition}

Die genaue Form von \(n(s,\phi)\) hängt vom gewählten mechanischen Modell ab.

\RESEARCH{Präzise Definition der Versatzrichtung anhand des Fahrzeugmodells.}
\OPEN{Konsistente Integration lokaler Gleisrahmen und globaler CRS-Transformationen in das Schwerpunktmodell.}

\section{Schwerpunktprinzip}

\begin{definition}[Schwerpunktprinzip]
Ein Eisenbahn-Alignment erfüllt das Schwerpunktprinzip, wenn die induzierte
Referenztrajektorie \(\gamma_c(s)\) vorgegebene Glattheits- und dynamische
Eigenschaften besitzt.
\end{definition}

In dieser Formulierung ist nicht das Gleis selbst, sondern die induzierte
Trajektorie des Bezugspunkts der primäre Entwurfsgegenstand. Die
Entwurfsaufgabe rückt damit näher an fahrzeugorientiertes Engineering als an
rein geometrische Konstruktion.

\section{Variationsformulierung}

\begin{definition}[Schwerpunktfunktional]
Es sei
\[
J_c[\gamma_c]=\int_0^L\|\gamma_c''(s)\|^2\,\dd s.
\]
\end{definition}

Dieses Funktional bestraft die Krümmung der Referenztrajektorie und fördert
damit glatte Bewegung.

\begin{definition}[Schwerpunkt-Optimierungsproblem]
Gesucht ist
\[
(\kappa,\phi),
\]
sodass die induzierte Trajektorie \(\gamma_c\) unter den pose3-Dynamiken das
Funktional
\[
J_c[\gamma_c]
\]
minimiert.
\end{definition}

Diese Formulierung erweitert die Variationsinterpretation des
Übergangsentwurfs von Krümmungsgesetzen allein auf fahrzeugrelevante
abgeleitete Bewegung.

\section{Kopplung an pose3}

Die Trajektorie \(\gamma_c\) hängt ab von
\[
\gamma(s),\quad\theta(s),\quad\phi(s).
\]
Damit wird das Schwerpunktproblem zu einem beschränkten Optimierungsproblem in
den Variablen
\[
\kappa(s),\quad\phi(s).
\]

Schwerpunkt-Trassierung liegt somit natürlich im zuvor eingeführten
pose3-basierten Zustandsmodell und in der weiter gefassten Sicht der
Trassierungsoptimierung
\cite{Kufver1997MathematicalDescription,Kufver2000RailwayAlignment}.

\section{Interpretation}

Die Schwerpunktformulierung verschiebt die Sicht
\begin{itemize}
\item von der Gleisgeometrie
\item zur Fahrzeugbewegung.
\end{itemize}

Es entsteht ein einheitlicher Rahmen, in dem Geometrie, Überhöhung und Dynamik
gleichzeitig optimiert werden. Diese Sicht entspricht der allgemeinen
Feststellung, dass Übergangs- und Alignment-Qualität nicht allein anhand der
Geometrie beurteilt werden kann
\cite{BrustadDalmo2020Review,EN13803_2017}.

\section{Bezug zum klassischen Entwurf}

\begin{proposition}
Klassische Übergangsbögen entsprechen Sonderfällen der
Schwerpunktformulierung, in denen die Referenztrajektorie mit der
Gleismittellinie zusammenfällt.
\end{proposition}

In diesem Fall reduziert sich das Problem auf krümmungsbasierte Optimierung in
der Ebene. Der klassische Alignment-Entwurf wird daher nicht widerlegt,
sondern als Sonderfall einer allgemeineren Formulierung eingebettet.

\section{Vorteile}
\begin{itemize}
\item einheitliche Modellierung von Geometrie und Dynamik,
\item natürliche Integration der Überhöhung,
\item Kompatibilität mit Optimierungsrahmen und
\item Erweiterbarkeit auf fahrzeugspezifische Modelle.
\end{itemize}

Diese Vorteile folgen daraus, dass das Modell Alignment als zustandsbestimmtes
physisches System und nicht als statische geometrische Linie behandelt.

\section{Offene Probleme}
\OPEN{Präzise mechanische Interpretation der Referenztrajektorie.}
\OPEN{Integration mit vollständigen Fahrdynamikmodellen.}
\OPEN{Erweiterung auf netzweiten Alignment-Entwurf.}
\OPEN{Ausdrückliche Behandlung von CRS, lokalen Gleisrahmen und Koordinatentransformationen in großräumigen Alignment-Systemen.}

\section{Verbindung zum AIM-Rahmen}
\begin{summary}
Schwerpunkt-Trassierung erweitert den klassischen Alignment-Entwurf
begrifflich und mathematisch.

Statt geometrische Elemente vorzuschreiben, definiert sie die Entwurfsaufgabe
anhand einer aus dem Eisenbahnzustand abgeleiteten Referenztrajektorie.

Dies entspricht der AIM-Sicht des Alignments als funktionalem Objekt, das von
Krümmung und deren abgeleiteten Größen bestimmt wird.
\end{summary}

\section{Fahrzeugmodell}

\subsection{Starrkörperapproximation}
Das Eisenbahnfahrzeug wird als Starrkörper mit charakteristischem
Massenschwerpunkt modelliert.

\begin{definition}[Massenschwerpunkt]
Es bezeichne
\[
C(s)\in\R^3
\]
die räumliche Position des Fahrzeugschwerpunkts.
\end{definition}

Für ebene Gleisgeometrie betrachten wir die Projektion
\[
\gamma_c(s)\in\R^2
\]
der Schwerpunkttrajektorie.

Dieses Modell ist bewusst vereinfacht. Es ersetzt keine vollständige
Mehrkörpersimulation, sondern führt eine mathematisch handhabbare
fahrzeugbezogene Zustandsebene ein.

\subsection{Gleiskoordinatenrahmen}
\begin{definition}[Frenet-Rahmen]
Es seien
\[
t(s)=(\cos\theta(s),\sin\theta(s))
\]
die Tangentenrichtung und
\[
n(s)=(-\sin\theta(s),\cos\theta(s))
\]
die Normalenrichtung.
\end{definition}

Das Paar \((t,n)\) definiert ein lokales, an das Gleis gebundenes
Koordinatensystem. Dieser lokale Rahmen genügt für die vorliegende
Formulierung; künftige Erweiterungen sollen ihn konsistent mit einem globalen
Projekt-CRS verbinden.

\subsection{Überhöhungsdrehung}
\begin{definition}[Gedrehte Normale]
Der Überhöhungswinkel \(\phi(s)\) definiert eine Drehung des Querschnitts.
\end{definition}

In einem vereinfachten Modell bleibt die laterale Versatzrichtung an
\(n(s)\) ausgerichtet, während vertikale Wirkungen implizit erfasst werden.
Diese Vereinfachung ist auf der vorliegenden begrifflichen Stufe vertretbar;
eine vollständige Behandlung muss den dreidimensional gedrehten Querschnitt
und dessen Transformation in den globalen Rahmen ausdrücklich einbeziehen.

\subsection{Versatz des Massenschwerpunkts}
\begin{definition}[Versatzmodell]
Es bezeichne \(d>0\) einen charakteristischen lateralen Versatz. Die
Schwerpunkttrajektorie wird angenähert durch
\[
\gamma_c(s)=\gamma(s)+d\,n(s).
\]
\end{definition}

Verfeinerte Modelle können eine Abhängigkeit von \(\phi(s)\) einbeziehen, etwa
\[
\gamma_c(s)=\gamma(s)+d\,\cos\phi(s)\,n(s).
\]
\RESEARCH{Einbeziehung vollständiger dreidimensionaler Fahrzeuggeometrie.}

\section{Dynamik des Massenschwerpunkts}

\subsection{Geschwindigkeit und Beschleunigung}
\begin{proposition}
Die Geschwindigkeit des Massenschwerpunkts ist
\[
\gamma_c'(s)=t(s)+d\,n'(s).
\]
\end{proposition}

Mit
\[
n'(s)=-\kappa(s)\,t(s)
\]
folgt
\[
\gamma_c'(s)=(1-d\kappa(s))\,t(s).
\]

\begin{proposition}
Die Beschleunigung des Massenschwerpunkts ist proportional zu
\[
\gamma_c''(s)
=
-d\kappa'(s)\,t(s)
+(1-d\kappa(s))\kappa(s)\,n(s).
\]
\end{proposition}

Somit beeinflussen Krümmung und ihre Ableitung die Schwerpunktbewegung. Darin
liegt ein zentraler begrifflicher Vorteil der Formulierung: Sie macht die
dynamische Relevanz geometrischer Größen höherer Ordnung ausdrücklich.

\subsection{Effektive Kräfte}
\begin{definition}[Querbeschleunigung]
Die Querbeschleunigung des Massenschwerpunkts ist
\[
a_c(s)=v^2(1-d\kappa(s))\kappa(s).
\]
\end{definition}

Mit Überhöhung ergibt sich die effektive Beschleunigung
\[
a_{\mathrm{eff}}(s)=v^2\kappa(s)-g\sin\phi(s).
\]

Der zweite Ausdruck spiegelt dieselbe Krümmungs--Überhöhungs-Kopplung, die in
Eisenbahnentwurfsnormen und dynamischer Alignment-Interpretation erscheint
\cite{Klein1998Trassierung,EN13803_2017}.

\section{Gekoppeltes Entwurfsproblem}

\begin{definition}[Fahrzeugbezogene Zielfunktion]
Definiere
\[
J_{\mathrm{veh}}[\kappa,\phi]
=
\int_0^L\|\gamma_c''(s)\|^2\,\dd s.
\]
\end{definition}
Dieses Funktional misst unmittelbar die Glattheit der Schwerpunktbewegung.

\begin{definition}[Dynamische Zielfunktion]
Definiere
\[
J_{\mathrm{dyn}}[\kappa,\phi]
=
\int_0^L\Bigl(v^2\kappa(s)-g\sin\phi(s)\Bigr)^2\,\dd s.
\]
\end{definition}

\begin{definition}[Kombiniertes Fahrzeugfunktional]
Eine kombinierte Zielfunktion ist
\[
J[\kappa,\phi]
=
\alpha J_{\mathrm{veh}}
+\beta J_{\mathrm{dyn}}
+\gamma\int_0^L(\kappa'(s))^2\,\dd s
+\delta\int_0^L(\phi'(s))^2\,\dd s.
\]
\end{definition}

Sie verbindet geometrische Glättung, dynamischen Ausgleich und
Regularisierung höherer Ordnung in einem Optimierungsrahmen.

\section{Schwerpunkt-Trassierung (formale Definition)}
\begin{definition}[Schwerpunkt-Alignment]
Ein Schwerpunkt-Alignment ist ein Paar
\[
(\kappa,\phi),
\]
das das Funktional
\[
J[\kappa,\phi]
\]
unter folgenden Bedingungen minimiert:
\begin{itemize}
\item pose3-Dynamik,
\item Randbedingungen und
\item Engineering-Beschränkungen.
\end{itemize}
\end{definition}

Damit wird Alignment-Entwurf zur Optimierung einer fahrzeugbezogenen
Zustandstrajektorie statt zur bloßen Komposition geometrischer Elemente.

\section{Interpretation}
In dieser Formulierung ist das Gleis nicht mehr der primäre Entwurfsgegenstand;
die Bewegung des Fahrzeugs wird zentral.

Der klassische Alignment-Entwurf ergibt sich als Sonderfall mit
\[
d=0\quad\text{und}\quad\phi(s)=0.
\]

\section{Diskussion}
Das dargestellte Modell ist bewusst vereinfacht. Es erfasst jedoch bereits:
\begin{itemize}
\item die Kopplung zwischen Krümmung und Überhöhung,
\item den Einfluss von Krümmungsableitungen und
\item die Beziehung zwischen Geometrie und Dynamik.
\end{itemize}

Sein Hauptbeitrag ist deshalb begrifflich und strukturell: Es definiert eine
mathematisch kohärente Verbindung zwischen Alignment-Geometrie und
fahrzeugorientiertem Entwurfsdenken.

\RESEARCH{Erweiterung auf vollständige Mehrkörper-Fahrzeugmodelle.}
\RESEARCH{Kalibrierung anhand gemessener Gleisdaten.}
\OPEN{Integration in praktische Entwurfsabläufe.}

\section{Verbindung zum AIM-Rahmen}
\begin{summary}
Die Verbindung von pose3 und einem fahrzeugbezogenen Modell stellt eine
strenge Grundlage für Schwerpunkt-Trassierung bereit.

Sie überführt Alignment-Entwurf in ein Problem der Optimierung der Bewegung
eines repräsentativen physischen Systems.

Dies unterstützt die zentrale AIM-Idee, dass Alignment nicht lediglich ein
geometrisches Objekt, sondern eine funktionale, durch Krümmung und ihre
physischen Folgen bestimmte Entität ist.
\end{summary}
